\documentclass[article]{jss}

\usepackage{amsmath,amssymb,amsfonts,bm,array}
\usepackage{float}
\usepackage{booktabs}
\usepackage{placeins}
\usepackage{needspace}
\usepackage[normalem]{ulem}
\hypersetup{hypertexnames=false}
\providecommand{\sn}[1]{\textcolor{blue}{#1}}
\RecustomVerbatimEnvironment{CodeOutput}{Verbatim}{fontsize=\small}

%% -- Author/title metadata ---------------------------------------------------
\author{Sinian Zhang\\University of Minnesota \And
        Jianfeng Wang\\University of Minnesota \And
        Thierry Chekouo\\University of Minnesota}
\Plainauthor{Sinian Zhang, Jianfeng Wang, Thierry Chekouo}

\title{\pkg{IntegMultiReg}: An \proglang{R} Package for Integrative
  \proglang{Bayesian} Multi-Regression of Multi-Platform Biomarkers}
\Plaintitle{IntegMultiReg: An R Package for Integrative Bayesian
  Multi-Regression of Multi-Platform Biomarkers}
\Shorttitle{\pkg{IntegMultiReg}}

\Abstract{
  The integration of high-dimensional molecular data measured on several
  platforms can sharpen the discovery of prognostic biomarkers, but it is
  routinely undermined by a loss of sample size: only a subset of subjects is
  typically measured on \emph{all} platforms, and complete-case analyses
  discard the rest. The \proglang{R} package \pkg{IntegMultiReg} implements the
  integrative multi-regression (IMR) model of \citet{chekouo2017}, which avoids
  this loss by partitioning subjects into the availability subgroups of a Venn
  diagram, fitting one regression per subgroup, and sharing information across
  subgroups through a Markov random field (MRF) prior on the
  variable-selection indicators combined with non-local priors on the
  regression coefficients. Whereas the original method targeted time-to-event
  outcomes, \pkg{IntegMultiReg} additionally supports continuous and binary
  outcomes through a unified interface. Posterior inference is
  carried out by a Markov chain Monte Carlo sampler written in \proglang{C} for
  efficiency, and the package provides standard \code{print}, \code{summary},
  \code{plot}, \code{coef} and \code{predict} methods, out-of-sample prediction
  for new subjects, explicit computational options, and cross-validated
  assessment of predictive accuracy. We
  describe the model, the package design, and illustrate both biomarker
  selection and the effect of borrowing strength across subgroups on a
  simulated multi-platform data set.
}
\Keywords{\proglang{Bayesian} variable selection, data integration, multi-omics,
  Markov random field, non-local prior, missing platforms, \proglang{R}}
\Plainkeywords{Bayesian variable selection, data integration, multi-omics,
  Markov random field, non-local prior, missing platforms, R}

\Address{
  Sinian Zhang, Jianfeng Wang and Thierry Chekouo\\
  Division of Biostatistics and Health Data Science\\
  University of Minnesota\\
  2221 University Ave SE, Suite 200\\
  Minneapolis, MN 55414, United States of America\\
  E-mails: \email{zhan9381@umn.edu}, \email{wan01188@umn.edu},
  \email{tchekouo@umn.edu}
}

\begin{document}

\clearpage

%% ===========================================================================
\section{Introduction}\label{sec:intro}
%% ===========================================================================

Modern biomedical studies characterise the same patients on multiple molecular
platforms -- for example gene expression, protein abundance, methylation or
metabolite levels -- in the hope that combining them improves the identification
of biomarkers that are associated with a clinical outcome
\citep{chekouo2015,chekouo2017}. A practical obstacle, observed across many
large cohorts, is that the platforms are rarely measured on exactly the same
subjects: a substantial fraction of patients are profiled on only a subset of
the available platforms. The most common remedy, restricting the analysis to
subjects with complete data on \emph{all} platforms, can discard more than half
of the available samples and, with it, much of the statistical power
\citep{chekouo2017}.

\citet{chekouo2017} proposed an integrative multi-regression (IMR) framework
that circumvents this loss. The key idea is to organise subjects into the
$2^{K}-1$ non-empty regions of a $K$-set Venn diagram according to which of the
$K$ platforms they have measured, to fit a separate regression model within
each region, and then to \emph{borrow strength} across these regions so that
biomarkers supported in one region inform selection in the others. Borrowing is
achieved through a Markov random field (MRF) prior on the binary
variable-selection indicators \citep{li2010,stingo2011,peterson2015}, while
sparsity and the elimination of spurious small effects are enforced by non-local
(product moment) priors on the regression coefficients
\citep{johnson2012,rossell2017}. The original work focused on accelerated
failure time (AFT) models for censored survival outcomes; the present package
extends the same machinery to continuous and binary responses.

A number of \proglang{R} packages address \proglang{Bayesian} variable selection
for high-dimensional and/or multivariate data. \pkg{spikeSlabGAM}
\citep{scheipl2011} provides spike-and-slab selection for structured additive
regression terms. \pkg{BayesSUR} \citep{zhao2021}
fits seemingly unrelated regressions with a choice of priors for the variable
selection indicators (including an MRF prior) and for the residual covariance
\citep{bottolo2021}; \pkg{mBvs} \citep{lee2017} performs multivariate
\proglang{Bayesian} variable selection exploiting dependence among outcomes.
\pkg{BayesSUR} and \pkg{mBvs} focus on dependence among multiple responses. Their
multivariate-response setting differs from the platform-availability problem
considered here. \pkg{IntegMultiReg} addresses a different and complementary problem:
a single outcome is regressed on predictors that come from several platforms,
and the platforms are observed on overlapping but distinct subsets of subjects.
The package therefore makes the missing-platform structure -- rather than the
outcome covariance -- the central modelling object.

Several multi-omics packages address related integration tasks. \pkg{MOFA}
\citep{argelaguet2018} learns unsupervised latent factors and can accommodate
incomplete entries, whereas \pkg{BIPnet} \citep{chekouo2023bipnet},
\pkg{SIDA} \citep{safo2022sida}, and \pkg{RandMVLearn}
\citep{safo2025rand} perform supervised association, classification, or
prediction through shared multiview representations. These methods are useful
for discovering variation or association across data views, but they do not
implement the IMR construction of separate outcome regressions for arbitrary
platform-availability subgroups with MRF-linked feature-selection indicators.
This availability-subgroup formulation is the distinguishing target of
\pkg{IntegMultiReg}.

These comparisons separate two sources of complexity: dependence among outcomes
and incomplete overlap among predictor platforms. Table~\ref{tab:software-scope}
places IMR alongside representative regression tools. Its closest predecessor
is the study-specific \proglang{C} implementation supplied with
\citet{chekouo2017}. \pkg{IntegMultiReg} builds on that method by providing
validated data objects, standard \proglang{R} methods, prediction for new
subjects and reproducible analysis workflows, while extending the outcome
families to continuous and binary responses.

\begin{table}[t!]
\centering
\small
\begin{tabular}{@{}>{\raggedright\arraybackslash}p{2.8cm}
                >{\raggedright\arraybackslash}p{3.3cm}
                >{\raggedright\arraybackslash}p{3.2cm}
                >{\raggedright\arraybackslash}p{4.4cm}@{}}
\toprule
Software or model & Primary modelling target & Treatment of incomplete
platforms & Relation to \pkg{IntegMultiReg} \\
\midrule
\citet{chekouo2017} supplement &
Single right-censored outcome with mRNA, miRNA and methylation predictors &
Availability subgroups are modelled directly through IMR &
Methodological basis; study-specific \proglang{C} workflow, survival outcome
only. \\
\pkg{BayesSUR} &
Multiple continuous responses in a seemingly unrelated regression &
Common sample and predictor matrices are assumed after preprocessing &
Shares Bayesian variable selection and MRF ideas, but targets outcome
covariance rather than platform availability. \\
\pkg{mBvs} &
Multivariate outcome regression with dependence among outcomes &
Common subjects and predictors are assumed &
Complementary multivariate-response setting; no explicit Venn-diagram
subgroup structure. \\
\pkg{glmnet} &
Penalised generalized linear and Cox models &
Requires a complete design matrix or externally imputed features &
Useful prediction baseline; no posterior inclusion probabilities or MRF
borrowing across availability subgroups. \\
\pkg{IntegMultiReg} &
One outcome with predictors from several partly observed platforms &
Subjects are partitioned into observed-platform subgroups and linked by an MRF
prior &
Reusable \proglang{R} implementation with survival, binary and continuous
outcomes, S3 methods, prediction and cross-validation. \\
\bottomrule
\end{tabular}
\caption{\label{tab:software-scope}Scope of \pkg{IntegMultiReg} relative to
closely related software and modelling workflows.}
\end{table}

The same distinction applies across software environments
(Table~\ref{tab:cross-language}). MOFA2 provides the Python
\pkg{mofapy2} implementation as well as an R interface; it learns latent
factors from multiple omics views, including partially overlapping samples
\citep{mofaDocs2026}. The Python \pkg{mvlearn} library includes CCA and
multiview extensions \citep{mvlearnDocs2026}. Similarity network fusion (SNF)
was implemented in MATLAB with R software distributed alongside the paper
\citep{wang2014snf}. In Julia, \pkg{MultivariateStats.jl} supplies CCA for two
matrices with aligned observations \citep{juliaCca2026}. These provide concrete
alternatives for factor discovery, correlated representations or network-based
integration. Their documented modelling targets differ from IMR's separate
outcome regressions for availability subgroups linked by feature-selection
priors. This is a comparison of scope, not a cross-package performance ranking
or a claim that related capabilities are absent from any language.

\begin{table}[t!]
\centering\small
\begin{tabular}{@{}>{\raggedright\arraybackslash}p{3.0cm}>{\raggedright\arraybackslash}p{2.0cm}>{\raggedright\arraybackslash}p{4.1cm}>{\raggedright\arraybackslash}p{4.1cm}@{}}
\toprule
Software & Environment & Related functionality & Distinction from IMR \\
\midrule
MOFA2 / \pkg{mofapy2} & R, Python &
Unsupervised multi-omics factors; missing entries omitted from the likelihood &
Describes shared variation rather than subgroup-specific outcome regression. \\
\pkg{mvlearn} & Python & CCA, multiview CCA and related embeddings &
Learns correlated representations; the documented CCA family is not the IMR selection model. \\
SNF & MATLAB, R & Fuses patient-similarity networks across data types &
Network integration rather than MRF-linked variable selection. \\
\pkg{MultivariateStats.jl} & Julia & CCA with a common observation count in both matrices &
A two-view correlation analysis, not regression over arbitrary availability subgroups. \\
\bottomrule
\end{tabular}
\caption{\label{tab:cross-language}Selected integration tools in R, Python, MATLAB and Julia, grouped by
modelling target rather than implementation language alone. Sources are the project documentation and the
SNF methods paper; documentation was consulted on September 21, 2026.}
\end{table}

The package is available through
\url{https://CRAN.R-project.org/package=IntegMultiReg}. This article describes
version 0.2.0; until that version reaches CRAN the accompanying replication
materials install it from
\url{https://github.com/h7nian/IntegMultiReg}, whose \code{main} branch carries
it. Readers reproducing the results below should check
\code{packageVersion("IntegMultiReg")} against 0.2.0, because the interface
changed at that version. The remainder of the paper follows the
analysis workflow from model specification to empirical assessment.
Section~\ref{sec:model} describes the IMR model and the priors that link the
regression models across subgroups. Section~\ref{sec:package} presents the
design of the package: the main fitting function and its arguments, the fitted
object and its methods, and the prediction and cross-validation utilities.
Section~\ref{sec:illustration} illustrates the package on a simulated
multi-platform data set, examining biomarker recovery and comparing prediction with and without
MRF borrowing. Section~\ref{sec:reproducibility} records the
replication materials and computational workflow. Section~\ref{sec:summary}
concludes.


%% ===========================================================================
\section{The integrative multi-regression model}\label{sec:model}
%% ===========================================================================

\subsection{Data structure and availability subgroups}\label{sec:subgroups}

Suppose $K$ platforms are available, and let $\bm{X}^{(k)}$ denote the
$p_k$ standardized features measured on platform $k$, $k=1,\dots,K$. Because a
subject may be measured on any subset of the platforms, the subjects partition
into at most $2^{K}-1$ disjoint \emph{availability subgroups}, one for each
non-empty region of the $K$-set Venn diagram. We index a subgroup by the binary
string $s=(s_K\cdots s_1)$, where $s_k=1$ if the subjects in that subgroup have
platform $k$ observed and $s_k=0$ otherwise; the right-most digit refers to the
first platform. For instance, with $K=3$ platforms the string \code{011}
denotes subjects observed on platforms~1 and~2 but not~3, \code{100}
denotes platform~3 only, \code{101} denotes platforms~1 and~3, and
\code{111} denotes all three platforms.  In schematic output we write the
platforms as \code{P1}, \code{P2}, \dots, \code{PK}; the $g$th feature on
platform \code{Pk} is denoted by \code{Pk:Fg} (or \code{PFkg} in compact
tables).  The package output also keeps the original platform and feature names
whenever they are supplied by the user.

For three platforms, Table~\ref{tab:bitstring-key} is a compact
Venn-diagram legend.  We use this table-like display, together with
\code{plot\_subgroup\_sizes()}, because it remains readable when some Venn
regions are empty or removed by the subgroup-size threshold.

\begin{table}[!ht]
\centering
\begin{tabular}{c c c c l}
\toprule
Bitstring & \code{P1} & \code{P2} & \code{P3} &
Availability subgroup \\
\midrule
\code{001} & yes & no  & no  & \code{P1} only \\
\code{010} & no  & yes & no  & \code{P2} only \\
\code{100} & no  & no  & yes & \code{P3} only \\
\code{011} & yes & yes & no  & \code{P1}+\code{P2} \\
\code{101} & yes & no  & yes & \code{P1}+\code{P3} \\
\code{110} & no  & yes & yes & \code{P2}+\code{P3} \\
\code{111} & yes & yes & yes & all platforms \\
\bottomrule
\end{tabular}
\caption{\label{tab:bitstring-key}Bitstring notation for a three-platform
availability structure.  The right-most digit corresponds to \code{P1}.}
\end{table}

Within subgroup $s$, let $n_s$ be the number of subjects and let
$\bm{y}^{(s)}$ be the (possibly latent) outcome vector. A subgroup is included
in the analysis only if $n_s$ exceeds a user-specified threshold, so that each
fitted regression is supported by enough data.

\subsection{Within-subgroup regression models}\label{sec:regression}

For each retained subgroup $s$, the standardized features of the platforms
present in $s$, collected in the design matrix $\bm{X}^{(s)}$, together with the
clinical covariates $\bm{C}^{(s)}$ (which are always included), enter a linear
model for a latent Gaussian response $\bm{y}^{\star(s)}$,
%
\begin{equation}\label{eq:linear}
  \bm{y}^{\star(s)} \;=\; \beta_0^{(s)}\bm{1}_{n_s}
  + \bm{C}^{(s)}\bm{\beta}_{C}^{(s)}
  + \bm{X}^{(s)}\bm{\beta}_{X}^{(s)} + \bm{\varepsilon}^{(s)},
  \qquad \bm{\varepsilon}^{(s)}\sim\mathcal{N}\!\left(\bm{0},\sigma_s^{2}\bm{I}_{n_s}\right).
\end{equation}
%
Here $\bm{\beta}_{C}^{(s)}$ and $\bm{\beta}_{X}^{(s)}$ are, respectively, the
coefficient vectors for the always-included clinical covariates and the
candidate molecular features. Depending on the outcome type,
$\bm{y}^{\star(s)}$ is defined as follows:
the notation $\bm{y}^{\star(s)}$ is a working response used to keep the
regression and variable-selection layers common across outcome types, not a new
observed endpoint shared by all analyses.  Once this working response is
observed or augmented, the same pMOM prior and MRF borrowing structure can be
used for feature selection.
%
\begin{itemize}
\item \emph{continuous} outcomes set $\bm{y}^{\star(s)}=\bm{y}^{(s)}$ directly;
\item \emph{right-censored} survival outcomes use an AFT specification in which
  $y^{\star(s)}_n=\log(t^{(s)}_n)$ for an event and, for a censored observation,
  $y^{\star(s)}_n>\log(t^{(s)}_n)$ is imputed by data augmentation
  \citep{tanner1987}, where $t^{(s)}_n$ denotes the observed event or censoring
  time for subject $n$ in subgroup $s$;
\item \emph{binary} outcomes use a probit data-augmentation scheme in which the
  latent $y^{\star(s)}_n$ is the underlying Gaussian utility whose sign gives the
  observed $0/1$ label \citep{albert1993}. This latent response is a
  computational augmentation, not an additional observed biological endpoint;
  the user still supplies only the observed binary labels, and the Gaussian
  utility is imputed inside the sampler.  Related Bayesian predictive
  imaging-genetics models use the same latent-Gaussian idea for binary disease
  classification
  \citep{chekouo2016aoas}.
\end{itemize}
%
A conjugate inverse-gamma prior with shape $\alpha$ and rate $\psi$ is placed on
the error variance $\sigma_s^2$. The always-included intercept and clinical
coefficients also have first-order pMOM priors, with scale $h_0$, as in
Section~3.1 of the original paper. A large $h_0$ makes this prior diffuse;
it does not make it a normal prior, and $h_0$ is not its marginal variance. Detailed likelihood and
data-augmentation forms for the continuous and binary extensions are given in
Appendices~\ref{sec:appendix-continuous} and~\ref{sec:appendix-binary}.

\subsection{Non-local priors for variable selection}\label{sec:nlp}

Selection is performed by binary indicators $\gamma^{(s)}_{kg}\in\{0,1\}$ that
flag whether feature $g$ of platform $k$ enters the regression for subgroup $s$.
Conditional on $\gamma^{(s)}_{kg}=0$ the coefficient is set to zero; conditional
on $\gamma^{(s)}_{kg}=1$ it follows a first-order non-local product-moment
(pMOM) prior \citep{johnson2012},
%
\begin{equation}\label{eq:pmom}
  p\!\left(\beta^{(s)}_{kg}\mid\gamma^{(s)}_{kg}=1\right)\;\propto\;
  \left(\beta^{(s)}_{kg}\right)^{2}
  \exp\!\left\{-\frac{1}{2\,\tau_k\,\sigma^2_s}\left(\beta^{(s)}_{kg}\right)^{2}\right\}.
\end{equation}
%
Because the density~\eqref{eq:pmom} vanishes at the origin, non-local priors
assign vanishing prior mass to negligible effects and thereby discard models
that contain unnecessary predictors, yielding more parsimonious and more
interpretable biomarker sets than the usual local (spike-and-slab) priors. The
scale $\tau_k$ is controlled by the argument \code{molecular\_prior\_scale}.

\subsection{A Markov random field prior to borrow strength}\label{sec:mrf}

To encourage the \emph{same} biomarkers to be selected across the regression
models, the vector of indicators
$\bm{\gamma}_{k,g}=(\gamma^{(s)}_{kg})_{s}$ for feature $g$ of platform $k$ is
given a Markov random field prior,
%
\begin{equation}\label{eq:mrf}
  p\!\left(\bm{\gamma}_{k,g}\mid\nu_k,\bm{\Theta}_k\right)
  \;=\; Z(\nu_k,\bm{\Theta}_k)^{-1}
  \exp\!\left(\nu_k\,\bm{1}^{\top}\bm{\gamma}_{k,g}
  + \bm{\gamma}_{k,g}^{\top}\bm{\Theta}_k\,\bm{\gamma}_{k,g}\right),
\end{equation}
%
where $\nu_k$ is a sparsity parameter (the prior log-odds of inclusion for
platform $k$) and the symmetric matrix $\bm{\Theta}_k$ has zero diagonal and
non-negative off-diagonal entries $\theta^{(k)}_{ss'}$ that measure how strongly
subgroups $s$ and $s'$ share biomarkers of platform $k$. A large
$\theta^{(k)}_{ss'}$ favours co-inclusion of the same feature in
the two subgroups, so a strong marker discovered in a large subgroup can be
recovered in a smaller one. The off-diagonal entries are assigned independent
gamma priors with shape $\alpha_0$ and rate $\beta_0$ (distinct from the
intercept $\beta_0^{(s)}$ of Equation~\ref{eq:linear}). Setting all
$\theta^{(k)}_{ss'}=0$ removes the dependence and recovers a non-integrative
\proglang{Bayesian} multi-step (BMS) analysis that fits each subgroup
separately; this variant is available for comparison through the argument
\code{method}.

\subsection{Posterior inference and prediction}\label{sec:inference}

The regression coefficients and variance parameters are integrated out using a
Laplace approximation to the marginal likelihood \citep{johnson2012}, and the
remaining parameters -- the selection indicators $\bm{\gamma}$, the MRF
interaction matrices $\bm{\Theta}_k$, and the augmented latent responses -- are
updated by a Metropolis-within-Gibbs Markov chain Monte Carlo (MCMC) sampler.
The post-burn-in draws yield, for every platform and subgroup, the marginal
posterior inclusion probability (mPIP) $\hat\gamma^{(s)}_{kg}$, which is the
package's primary measure of biomarker importance.

Predictions for new subjects are obtained by \proglang{Bayesian} model
averaging over the sampled selection models, weighted by their marginal
likelihoods. A new subject is first assigned to the subgroup matching its own
pattern of observed platforms; its features are standardized using the training
centring and scaling factors, and the averaged regression yields a predicted
latent response (a survival/continuous value, or a probability through the
probit link for binary outcomes).
Concretely, for a new subject in subgroup $s$, the implementation keeps the
highest-posterior distinct selection models and uses the following
model average:
\[
  \widehat{y}^{\star(s)}_{\mathrm{new}}
  =
  \sum_{\ell=1}^{L}
  w_{\ell}^{(s)}
  \bm{z}_{\mathrm{new}}\!\left(\bm{\gamma}_{\ell}^{(s)}\right)^{\top}
  \widehat{\bm{\beta}}_{\ell}^{(s)}, \qquad
  w_{\ell}^{(s)}
  =
  \frac{\exp\{\log p(\bm{\gamma}_{\ell}^{(s)},\bm{\Theta}\mid
  \widehat{\bm{y}}^{\star})\}}
  {\sum_{r=1}^{L}\exp\{\log p(\bm{\gamma}_{r}^{(s)},\bm{\Theta}\mid
  \widehat{\bm{y}}^{\star})\}},
\]
which is the package analogue of the model-averaged prediction formula in
\citet{chekouo2017}.  Here $\bm{z}_{\mathrm{new}}(\bm{\gamma}_{\ell}^{(s)})$
contains the intercept, clinical covariates and selected platform features for
selection model $\ell$.

\subsection{Analysis strategies represented by the package}\label{sec:strategies}

It is useful to distinguish three analysis strategies that are often conflated
in multi-platform studies.  A \emph{complete-case} strategy first restricts the
data to subjects measured on every platform and then fits a single regression
model.  This gives a simple design matrix, but it discards all subjects in the
platform-specific regions of the Venn diagram.  In the KIRC case study of
\citet{chekouo2017}, the complete three-platform region contains only a
minority of the subjects with at least one platform measured, so the
complete-case analysis is both smaller and scientifically narrower.

The \emph{Bayesian multi-step} (BMS) strategy keeps the availability subgroups
but removes the MRF interaction terms by setting the subgroup-dependence
parameters to zero.  Each subgroup then has its own variable-selection problem,
with the same non-local prior on active regression effects but without borrowing
strength across subgroups.  This is the no-borrowing comparator used by
\citet{chekouo2017} to isolate the contribution of the MRF prior.  BMS is
therefore a useful internal comparator: it answers the question of how much is
gained by linking the subgroup-specific selection indicators rather than simply
fitting all retained subgroups independently.

The \emph{integrative multi-regression} (IMR) strategy keeps the BMS
within-subgroup likelihoods and non-local priors, but adds the MRF prior in
Equation~\ref{eq:mrf}.  IMR is not an imputation method for missing molecular
features; no unobserved platform values are filled in.  Instead, it uses the
observed data in each availability subgroup and shares selection information at
the biomarker level whenever the same platform appears in multiple subgroups.
This distinction is important for interpretation.  A feature with high mPIP in
a small subgroup is supported by that subgroup's observed outcomes and
covariates, but its posterior probability may be stabilized by evidence for the
same feature in larger neighbouring subgroups.  Conversely, when a subgroup is
large enough to identify its own signal, the MRF prior need not force agreement;
the likelihood can still dominate the borrowing term.

In \pkg{IntegMultiReg}, the two subgroup-preserving strategies are selected by
\code{method = "imr"} and \code{method = "bms"}.  A complete-case analysis can
be produced by supplying only the rows with all platforms observed, but it is
not the default because the central use case of the package is precisely to
avoid unnecessary loss of partially observed subjects.

%% ===========================================================================
\section[The package IntegMultiReg]{The package \pkg{IntegMultiReg}}\label{sec:package}
%% ===========================================================================

The model determines what information is shared across subgroups; the package
interface determines how an analyst supplies, inspects and validates that
analysis. We first describe the data objects, then fit a model and examine
its summaries, uncertainty and predictions.

\pkg{IntegMultiReg} is written in \proglang{R} \citep{rcore} with the
computationally intensive MCMC sampler implemented in \proglang{C} and linked
against the GNU Scientific Library \citep[GSL,][]{galassi2009}, which must be
available on the system (\code{brew install gsl} on macOS,
\code{apt-get install libgsl-dev} on Debian/Ubuntu, or Rtools on Windows).

\subsection{Data workflow and object design}\label{sec:workflow}

The package provides a dedicated \code{"imr\_data"} class so that alignment is
validated independently of model fitting.  The constructor \code{imr\_data()}
checks identifiers, feature types, outcome coding and dimensions, records the
observed-platform pattern for every analysis-eligible subject, reports subjects
excluded for missing required outcome/covariate rows, and supports \code{print()},
\code{summary()} and \code{validate\_imr\_data()}.  The same object can be
passed to \code{predict()} when it contains new platform measurements and
covariates.  This avoids repeating the alignment logic separately for fitting
and prediction.

The package accepts ordinary rectangular \proglang{R} data frames and validates
them through the dedicated \code{"imr\_data"} container.  Each
platform is supplied as one data frame with an \code{id} column and numeric
features; the outcome and optional clinical covariates use the same identifier.
The fitting function then constructs the modelled data structure in five steps.
First, subjects are restricted to those with the required outcome (and
covariate) information and at least one molecular platform.  Second, a platform
presence matrix is created and converted to the bitstring convention described
in Section~\ref{sec:subgroups}.  Third, subgroups with at most \code{min\_subgroup\_size}
subjects are removed.  Fourth, outcome, covariate and platform rows are aligned
by \code{id} within every retained subgroup; this alignment is explicit, so the
input data frames need not share the same row order.  Finally, the molecular
features and clinical covariates are standardized within each subgroup, and the
corresponding centring and scaling factors are stored for prediction.

This preprocessing layer is implemented in \proglang{R}, where the inputs can be
validated with informative error messages.  The computationally intensive
updates are delegated to compiled \code{.Call} routines, while the
\proglang{R} wrapper records the call, hyperparameters, platform labels,
feature names, subgroup sizes and normalization constants needed by downstream
methods.

\begin{table}[!ht]
\centering
\begin{tabular}{@{}>{\raggedright\arraybackslash}p{3.7cm}
                >{\raggedright\arraybackslash}p{10.5cm}@{}}
\toprule
Component & Role in the fitted \code{"imr"} object \\
\midrule
\code{control} & Outcome and response scale, priors, MCMC settings, seed, sampling and numerical conventions. \\
\code{model} & Platform and feature names, covariates, availability subgroups and their mappings. \\
\code{preprocessing} & Validated raw inputs, transformed matrices, centring/scaling factors and formula terms. \\
\code{posterior} & Inclusion probabilities, interaction draws and means, selection draws, latent-response means and log-posterior trace. \\
\bottomrule
\end{tabular}
\caption{\label{tab:object}Main components of the fitted \code{"imr"} object.}
\end{table}

Two design choices are particularly important in applied multi-omics analyses.
The first is that clinical covariates are carried through the same subgroup
alignment and standardization workflow as molecular features, but are always
included rather than selected.  This matches the case-study use in which age,
sex, stage and grade are adjustment variables rather than candidate
biomarkers.  The second is that the prediction code repeats the training
availability logic for new subjects.  A user may supply only the platforms that
are available in the test set; each new subject is routed to a retained
availability subgroup, standardized with the training means and standard
deviations for that subgroup, and predicted by the corresponding averaged
regression model.  Subjects whose observed-platform pattern was not retained
during fitting are reported and omitted rather than silently assigned to an
incompatible model.

\subsection{Fitting the model}\label{sec:fit}

The main function is \code{imr()}. Its arguments are deliberately
\emph{orthogonal}: each controls a single, distinct aspect of the analysis, so
that data, likelihood, model, prior, computation and output can be varied
independently. Table~\ref{tab:args} groups the arguments by their role, in the
grouping used in the discussion below.

\begin{table}[!htbp]
\centering\small
\begin{tabular}{@{}>{\raggedright\arraybackslash}p{2.5cm}
                >{\raggedright\arraybackslash}p{4.6cm}
                >{\raggedright\arraybackslash}p{6.8cm}@{}}
\toprule
Role & Argument & Description \\
\midrule
\emph{Data}        & \code{x} & List of per-platform data frames; each has an \code{id} column and feature columns. \\
                   & \code{outcome}          & Data frame with \code{id} and the response (and event/censoring columns for survival). \\
                   & \code{covariates}              & Optional clinical covariates (always included, never selected). \\
\emph{Likelihood}  & \code{outcome\_type}     & \code{"right.censored"}, \code{"binary"} or \code{"continuous"}. \\
\emph{Model}       & \code{method}           & \code{"imr"} (integrative) or \code{"bms"} (independent subgroups). \\
\emph{Filtering}   & \code{min\_subgroup\_size}            & Minimum subgroup size; smaller subgroups are dropped. \\
\emph{Prior}       & \code{nu}               & Prior log-odds of inclusion, one per platform (sparsity). \\
                   & \code{molecular\_prior\_scale}               & Scale of the non-local prior on the regression effects. \\
                   & \code{forced\_prior\_scale}               & pMOM scale for intercept and clinical effects. \\
                   & \code{residual\_prior}    & Shape and rate of the error-variance prior. \\
                   & \code{interaction\_prior}    & Shape and rate of the MRF interaction prior. \\
\emph{Computation} & \code{draws}, \code{burnin}      & Retained and discarded MCMC iterations, respectively. \\
                   & \code{seed}             & Optional seed; drawn from R's RNG when \code{NULL} (default). \\
                   & \code{sampler\_method} & Legacy updates or the corrected paper-target sampler. \\
                   & \code{initial} & Optional named selection and IMR interaction matrices for explicit chain starts. \\
                   & \code{prior\_indexing} & Standard precision blocks or the released-code indexing convention. \\
                   & \code{laplace\_max\_iter}, \code{laplace\_tolerance} & Stage-specific coefficient-mode iteration limits and convergence tolerance. \\
\emph{Preprocessing} & \code{standardize} & Standardize predictors by default; disable for already transformed inputs. \\
                   & \code{survival\_scale} & Log positive survival times by default; identity is an explicit historical alternative. \\
\emph{Output}      & \code{verbose}          & Print sampler progress (default \code{FALSE}). \\
\bottomrule
\end{tabular}
\caption{\label{tab:args}Arguments of \code{imr()}, grouped by the orthogonal
aspect of the analysis that each one controls.}
\end{table}

A typical call, on the simulated data set \code{simIMR} shipped with the
package (three platforms -- \code{genomic}, \code{proteomic},
\code{metabolomic} -- with 20, 10 and 8 features measured on overlapping subsets
of 300 subjects, plus three clinical covariates), is

\begin{CodeChunk}
\begin{CodeInput}
R> options(width = 72)
R> library("IntegMultiReg")
R> data("simIMR")
R> analysis_data <- imr_data(
+    platforms    = simIMR$platforms,
+    outcome      = simIMR$outcome.survival,
+    covariates   = simIMR$covariates,
+    outcome_type = "right.censored")
R> fit <- imr(
+    x = analysis_data,
+    nu                 = c(-4, -3, -4),
+    draws = 8000, burnin = 2000,
+    sampler_method = "paper",
+    min_subgroup_size              = 30,
+    seed               = 1)
\end{CodeInput}
\end{CodeChunk}

The original component-wise call remains supported: platform matrices may
still be supplied as a list rather than through \code{imr\_data()}. The
argument names changed at 0.2.0, however, and the pre-0.2.0 aliases are not
accepted; \code{inst/MIGRATION.md} maps the old names to the new ones.
For analyses organised as one subject-level data frame, a formula method is
also available.  A binary response and adjustment covariates can be specified
by the formula \code{y \textasciitilde{} age + sex + stage}, with the subject
table supplied through \code{data = clinical} and the assay tables through
\code{platforms = assays}. Standard \code{model.frame()} and
\code{model.matrix()} processing then handles numeric and factor covariates,
while subject identifiers continue to align the separate platform tables.
Prediction from raw clinical data reuses the training terms, factor levels,
contrasts and supported transformation bases. The identifier is excluded when
expanding a dot formula. Because the IMR model includes an intercept and does
not implement offsets, formulas that remove the intercept or specify an offset
are rejected explicitly.

The function returns an object of class \code{"imr"}. Printing it gives a
compact overview: the outcome type, the method, the modelled availability
subgroups with their sizes, and the number of features selected per platform at
a marginal inclusion probability above the user-specified
reporting threshold, whose default is $0.5$.  The threshold is an explicit
argument of \code{print()} and \code{summary()}, as shown below.
The high-dimensional posterior summaries remain accessible as ordinary
\proglang{R} objects: for example, \code{coef(fit)\$genomic} is a matrix whose
rows are availability subgroups and whose columns are genomic features.

The \code{print()} method also gives a platform key so that the bitstrings
can be read without consulting a separate Venn diagram.  When a compact ranked
feature list is desired, \code{print(fit, rank = TRUE, top = 3)} appends the
top features per platform, ordered by each feature's maximum mPIP over
availability subgroups.

\begin{CodeChunk}
\begin{CodeInput}
R> print(fit, threshold = 0.5)
\end{CodeInput}
\begin{CodeOutput}
Integrative Bayesian Multi-Platform Regression (IMR)
----------------------------------------------------
Call:
  imr.imr_data(x = analysis_data, nu = c(-4, -3, -4), draws = 8000, 
    burnin = 2000, sampler_method = "paper", min_subgroup_size = 30, 
    seed = 1)

Outcome type : right.censored
Method       : IMR
Platforms    : 3 (genomic, proteomic, metabolomic)
MCMC         : 8000 retained draws after 2000 burn-in

Platform key:
  P1 = genomic
  P2 = proteomic
  P3 = metabolomic

Availability subgroups modelled (bitstring : platforms : size):
  011 : P1 + P2      : 120
  100 : P3           : 60
  101 : P1 + P3      : 60
  111 : P1 + P2 + P3 : 60

Features with mPIP > 0.50 (in any subgroup):
  genomic      : 5 of 20
  proteomic    : 2 of 10
  metabolomic  : 2 of 8
\end{CodeOutput}
\end{CodeChunk}

Setting \code{rank = TRUE} appends a short per-platform ranking of the
highest-mPIP features to the same overview, with \code{top} controlling how many
are listed:

\begin{CodeChunk}
\begin{CodeInput}
R> print(fit, rank = TRUE, top = 3)
\end{CodeInput}
\begin{CodeOutput}
Integrative Bayesian Multi-Platform Regression (IMR)
----------------------------------------------------
Call:
  imr.imr_data(x = analysis_data, nu = c(-4, -3, -4), draws = 8000, 
    burnin = 2000, sampler_method = "paper", min_subgroup_size = 30, 
    seed = 1)

Outcome type : right.censored
Method       : IMR
Platforms    : 3 (genomic, proteomic, metabolomic)
MCMC         : 8000 retained draws after 2000 burn-in

Platform key:
  P1 = genomic
  P2 = proteomic
  P3 = metabolomic

Availability subgroups modelled (bitstring : platforms : size):
  011 : P1 + P2      : 120
  100 : P3           : 60
  101 : P1 + P3      : 60
  111 : P1 + P2 + P3 : 60

Features with mPIP > 0.50 (in any subgroup):
  genomic      : 5 of 20
  proteomic    : 2 of 10
  metabolomic  : 2 of 8

Top 3 ranked features by maximum subgroup mPIP:
  genomic
     feature max_mpip subgroup
         G01        1      011
         G02        1      011
         G03        1      011

  proteomic
     feature max_mpip subgroup
         P01    1.000      011
         P02    1.000      011
         P06    0.178      111

  metabolomic
     feature max_mpip subgroup
         M01    1.000      100
         M02    0.505      111
         M06    0.403      111

\end{CodeOutput}
\end{CodeChunk}

\subsection{Choosing analysis settings}\label{sec:settings}

The preceding example specifies a complete analysis, but its settings should
be reconsidered for each study.  The subgroup-size threshold \code{min\_subgroup\_size} is the first such choice.  A
small value retains more availability patterns, but may create regressions with
too few events or too little outcome variation to support stable inference.  A
large value removes sparse regions of the Venn diagram and reduces computation,
but may discard scientifically relevant platform combinations.  In practice we
recommend first inspecting the platform-availability table, then setting
\code{min\_subgroup\_size} so that every retained subgroup has enough outcome information for
the chosen likelihood.  The function \code{plot\_subgroup\_sizes()} provides a
quick visual check after fitting.

The sparsity parameter \code{nu} is platform-specific because different assays
often have different effective dimensions and signal densities.  More negative
values favour smaller selected sets.  The KIRC analysis follows
\citet{chekouo2017} and uses \code{nu = c(-4, -3, -4)} for mRNA, miRNA and
methylation; in the simulated examples the same pattern is used for consistency.
The non-local prior scale \code{molecular\_prior\_scale} controls how strongly active effects are
kept away from zero.  Larger values allow broader slab effects, whereas smaller
values place active effects closer to zero but still retain the pMOM density's
zero-at-the-origin behaviour.  The MRF hyperparameters in
\code{interaction\_prior} govern the prior strength of borrowing; setting
\code{method = "bms"} is the most direct sensitivity analysis for this component
because it removes the MRF interactions while leaving the rest of the modelling
pipeline unchanged.

These prior choices define the target model. The sampler convention must also
be stated, because the published model and the released implementation differ
in their updates.

The stated symmetric MRF has conditional log-odds
$\nu_k + 2\sum_{s'\ne s}\theta^{(k)}_{ss'}\gamma^{(s')}_{kg}$.
The released selection code uses a factor of one and omits the Hastings
correction when flip proposals cross an empty/full-model boundary.
The explicit \code{sampler\_method = "paper"} option implements the stated
conditional, the proposal correction and the negative Gamma rate term in the
log-posterior trace. The default \code{"legacy"} retains historical updates;
its invariant distribution is not the stated paper posterior. The main fit and predictive simulation comparison request \code{"paper"}
explicitly. The separate covariate-formula workflow retains the legacy sampler
for both candidates; its results are not combined with this comparison.
\code{prior\_indexing = "code2017"} and stage-specific
\code{laplace\_max\_iter} expose two further released-code conventions.
\code{standardize = FALSE} accepts externally transformed historical inputs;
users remain responsible for the scope of that external transformation.

Several package arguments are thresholds or size controls, but they affect
different parts of the workflow and should not be interpreted as the same
quantity.  Table~\ref{tab:thresholds} separates these controls.  In particular,
\code{min\_subgroup\_size} determines which availability subgroups are fitted, whereas the
\code{threshold} argument in \code{print()} and \code{summary()} only determines
which fitted features are reported as selected.  The number of candidate
features is fixed by the columns supplied in each platform data frame after any
external screening; \pkg{IntegMultiReg} then estimates their posterior
inclusion probabilities rather than choosing the input feature dimension
automatically.

\begin{table}[!htbp]
\centering
\small
\begin{tabular}{@{}>{\raggedright\arraybackslash}p{3.2cm}
                >{\raggedright\arraybackslash}p{4.4cm}
                >{\raggedright\arraybackslash}p{6.0cm}@{}}
\toprule
Control & Where used & Interpretation \\
\midrule
\code{min\_subgroup\_size} & \code{imr()} &
Minimum subgroup size for fitting; subgroups with at most this many
subjects are removed before MCMC. \\
\code{nu} & \code{imr()} &
Platform-specific prior log-odds of feature inclusion; more negative values
encode stronger sparsity. \\
\code{threshold} & \code{print()}, \code{summary()} &
Reporting cutoff for mPIP; it changes the displayed selected-feature list,
not the MCMC fit. \\
\code{rank}, \code{top} & \code{print()},
\code{plot\_top\_features()} &
\code{rank = TRUE} adds top-feature rankings to \code{print()}; \code{top}
sets how many highest-mPIP features are shown in the print output or ranked
feature plot. \\
\code{max\_models} & \code{predict()}, \code{cv\_imr()} &
Maximum number of distinct sampled selection models retained for
Bayesian model averaging in prediction. \\
\bottomrule
\end{tabular}
\caption{\label{tab:thresholds}Thresholds and size controls used by
\pkg{IntegMultiReg}.}
\end{table}

Finally, MCMC length should be scaled to the purpose of the run.  Short chains
are useful for checking data alignment, subgroup retention and the absence of
obvious numerical problems.  Longer chains should be used for reported mPIPs,
theta summaries and cross-validated prediction.  The package stores the full
log-posterior trace, and \code{plot(fit, type = "trace")} gives a quick
diagnostic display; for high-dimensional applications, we also recommend
repeating the final analysis under several seeds and checking that the leading
selected features are stable.


\subsection{The fitted object and its methods}\label{sec:methods}

The fitted \code{"imr"} object carries the posterior summaries together with
the metadata needed for prediction. Standard \proglang{S}3 methods and
diagnostic helpers are provided
(Table~\ref{tab:methods}). The \code{summary} method ranks, per platform, the
features whose marginal inclusion probability exceeds a threshold, together with
the subgroup in which that maximum is attained.


\Needspace{6\baselineskip}
The printed and summarized output therefore has two layers of names:
availability subgroups are labelled by the bitstrings in
Table~\ref{tab:bitstring-key}, while platforms and features use the names
stored in the fitted object.  In simulations those names are \code{genomic},
\code{proteomic}, \code{metabolomic} and feature labels such as \code{G01} or
\code{P01}; in generic notation they correspond to \code{P1}, \code{P2} and
features such as \code{P1:F01}.

\begin{table}[!htbp]
\centering
\begin{tabular}{l p{9.6cm}}
\toprule
Method & Description \\
\midrule
\code{print}   & Overview of the fit, platform key, availability bitstrings and selected-feature counts at a chosen mPIP threshold. \\
\code{summary} & Per-platform table of selected biomarkers and their inclusion probabilities. \\
\code{coef}    & List of per-platform matrices of marginal posterior inclusion probabilities (mPIP). \\
\code{plot}    & Heatmaps of the mPIP or MRF interactions and trace plots for the log-posterior, theta interactions or individual selection indicators. \\
\code{predict} & Predicted outcomes for new subjects by Bayesian model averaging. \\
\code{confint} & Equal-tail credible intervals from retained selection-indicator and theta draws. \\
\code{validate\_imr} & Structural and numerical validation of a saved fitted object. \\
\code{compare\_imr} & Descriptive comparison of several fits at a common mPIP threshold. \\
\bottomrule
\end{tabular}
\caption{\label{tab:methods}Methods and diagnostic helpers for objects of class
\code{"imr"}.}
\end{table}

\Needspace{30\baselineskip}
\begin{CodeChunk}
\begin{CodeInput}
R> summary(fit, threshold = 0.5)
\end{CodeInput}
\begin{CodeOutput}
Integrative Bayesian Multi-Platform Regression (IMR) -- summary
--------------------------------------------------------------
Outcome type : right.censored   Method: IMR
Selection threshold (mPIP) : 0.50

Platform 'genomic': 5 selected feature(s)
 feature max_mpip subgroup
     G01    1.000      011
     G02    1.000      011
     G03    1.000      011
     G19    0.712      011
     G05    0.630      101

Platform 'proteomic': 2 selected feature(s)
 feature max_mpip subgroup
     P01        1      011
     P02        1      011

Platform 'metabolomic': 2 selected feature(s)
 feature max_mpip subgroup
     M01    1.000      100
     M02    0.505      111

\end{CodeOutput}
\end{CodeChunk}

The printed ranking compresses each feature's subgroup-specific probabilities
to one maximum. The full matrices returned by \code{coef(fit)} retain these
differences, which the heatmap displays directly (Figure~\ref{fig:selection}).
Similar high-probability columns across subgroups indicate agreement in
selection, but this plot alone cannot attribute that agreement to the MRF
prior. The ranked view summarises the strongest feature scores
(Figure~\ref{fig:top}); the trace instead describes the sampled log-posterior
trajectory (Figure~\ref{fig:trace}). The IMR--BMS comparison in
Section~\ref{sec:gain} examines the consequences of including MRF borrowing.
Posterior standard deviations and credible intervals from the retained
selection-indicator and MRF-interaction draws are returned by
\code{posterior\_summary(fit)} or the standard \code{confint(fit)} method.
Parameter-specific traces are available through
\code{plot(fit, type = "theta\_trace", ...)} and
\code{plot(fit, type = "selection\_trace", ...)}.


\FloatBarrier

\begin{figure}[!htbp]
\centering
\includegraphics[width=\textwidth]{figures/selection.pdf}
\caption{\label{fig:selection}Marginal posterior inclusion probabilities
(mPIP) for the simulated right-censored example (8000 retained draws after
2000 burn-in iterations). Rows identify availability subgroups and columns
identify features. A shared grayscale maps mPIP from 0 (light) to 1 (dark),
allowing direct comparison across platforms. The display is produced by
\code{plot(fit, type = "selection")}.}
\end{figure}

\begin{figure}[!htbp]
\centering
\includegraphics[width=0.7\textwidth]{figures/top_features.pdf}
\caption{\label{fig:top}The ten largest platform-feature scores produced by
\code{plot\_top\_features(fit)}.  For each platform-feature pair, the score is
the highest mPIP attained among the availability subgroups containing that
platform, and each bar is annotated with the availability subgroup (bitstring)
in which that maximum was attained. Bars are coloured by platform and the dashed
line marks the $0.5$ reporting cutoff. Scores come from the same fit as
Figure~\ref{fig:selection}; a maximum does not imply high inclusion probability
in every subgroup.}
\end{figure}

\begin{figure}[!htbp]
\centering
\includegraphics[width=0.7\textwidth]{figures/trace.pdf}
\caption{\label{fig:trace}Log-posterior trajectory for the fit in Figure~\ref{fig:selection},
produced by \code{plot(fit, type = "trace")}. The dashed line separates the
2000 burn-in iterations from the 8000 retained iterations. This single-chain
display can reveal gross instability but does not establish convergence.}
\end{figure}


\FloatBarrier

\subsection{Coefficient and predictive intervals}

Inclusion probabilities describe feature selection; coefficient and predictive
intervals address uncertainty in effect sizes and outcomes. The optional
\code{posterior\_draws(fit)} stage
resamples whole retained selection states, preserving their empirical joint
subgroup distribution, and draws coefficients and residual variances
conditionally on each subgroup model. All active coefficients, including the
intercept and clinical effects, use the original first-order pMOM prior.
Inactive molecular coefficients remain exactly zero. Binary and censored
outcomes are augmented within these conditional chains using the observed data.

The function \code{posterior\_draws()} returns an \code{"imr\_posterior"}
object. Its \code{confint()} method gives coefficient intervals on the
transformed predictor scales. The method
\code{predict()} on this object with \code{type = "mean"} gives intervals for conditional
response means or binary event probabilities; \code{type = "response"}
adds future outcome variability. For log-time survival fits these are
transformed to time units and the point summaries are medians. A posterior
mean on the time scale need not exist under an inverse-gamma variance mixture.
Future censoring times are not generated.

These intervals retain the original sampler's Laplace approximation to model
weights. Conditional split R-hat is reported for each sampled model, with an
alert above 1.05 or when undefined; it does not diagnose the selection chain.
Both stages require adequate exploration. These additions do not introduce
clinical covariate selection.

\subsection{Prediction and cross-validation}\label{sec:predict}

The \code{predict} method produces predictions for new subjects, automatically
routing each subject to the subgroup that matches its observed platforms and
applying the stored training transformation. The standard \code{predict()}
generic handles both point predictions and conditional posterior summaries.

\Needspace{18\baselineskip}
\begin{CodeChunk}
\begin{CodeInput}
R> newX <- simIMR$platforms$genomic[1:20, ]
R> newP <- simIMR$platforms$proteomic[1:20, ]
R> pred <- predict(fit, newdata = list(newX, newP),
+    covariates = simIMR$covariates)
R> print(pred[["model:011"]][1:3, ], digits = 3)
\end{CodeInput}
\begin{CodeOutput}
  id prediction
1  1      0.397
2  2     -0.228
3  3      0.290
\end{CodeOutput}
\end{CodeChunk}

For the default log-time survival fit these point predictions are on the log-time
scale. The returned predictions retain full floating-point precision; rounding in the
display above is controlled by the user through the printing method and does
not alter the returned values.

Prediction for supplied subjects and validation of the fitting procedure answer
different questions. For the latter, \code{cv\_imr(cv\_method = "refit")}
recomputes preprocessing, formula transformations, selection draws and
predictive weights within each training fold. This is the workflow used in
the manuscript's predictive comparison. Its cost is approximately one full fit
per fold and round; binary predictions average model-specific probabilities
on the response scale.

The two post-fit modes instead reuse a model fitted to the full data. The
default \code{cv\_method = "legacy"} uses ranked distinct selection states
and historical scoring. The \code{"importance"} mode preserves all retained
states, including their multiplicities, for the importance average in
Equations (6)--(7) of \citet{chekouo2017}. Both use full-fit transformed
predictors and augmented response means; the response plug-in is stated in
Section 4.1 of that article. Held-out outcomes enter inverse-density weights
as part of the importance correction. These workflows therefore require their
own interpretation and should not be labelled training-fold refits.

Mode names supply defaults. The optional arguments \code{ridge},
\code{model\_set}, \code{df\_method} and \code{score\_method} choose the
coefficient penalty, selection-state collection, predictive degrees of freedom
and metric convention independently. Paper-oriented post-fit CV uses
\code{ridge = 0}, \code{model\_set = "draws"} and
\code{df\_method = "fractional"}. The released code instead uses
\code{model\_set = "ranked\_unique"}, \code{max\_models = 100},
\code{df\_method = "legacy\_integer"} and \code{score\_method = "legacy"},
also with zero ridge. The package retains 0.001 as its default penalty.

These numerical choices also determine when prediction is well defined.
A singular unpenalized solve stops with round, fold and subgroup context.
This restriction is material for the bundled KIRC inputs: in a 20,000-state
paper-sampler pilot with clinical covariates, subgroup 011 had more selected
columns (including intercept and covariates) than training subjects in
22.065\% to 28.060\% of retained states, depending on the fold. An independently
reconstructed design had 64 rows, 69 columns and rank 64. Thus this pilot does
not supply an unpenalized all-state CV estimate. A positive-penalty control
completed, but changes the estimator and is not a paper-reference result.
These pilot observations do not replace the original-scale experiments.


Once a validation workflow has been chosen, its predictions can be summarised
in two ways. The returned fields \code{pooled} and \code{fold\_mean} distinguish pooled
out-of-fold accuracy from the average fold-specific score. The ordinary
\code{predictions} data frame records subjects, subgroups, rounds, folds and
predictions; \code{control} records effective settings and a reusable fold
table. Supplying \code{folds = cv\$control\$folds} pairs analyses by actual
membership and preserves post-fit row order. Identical seeds do not pair the
R refit and GSL post-fit generators. New fits can request
\code{fold\_rng = "continue"} to continue the fit's GSL stream, matching the
released program's convention. This does not assert that preceding random
consumption or initial states match that program.

The preceding choices concern the analysis itself. To distribute its computation,
all modes support \code{workers} for deterministic process parallelism.
Settings and partitions are resolved before dispatch; each worker needs its own
fit and workspace. Hyperparameter or feature-screening choices selected from
data require outer validation. Existing results retain their interpretation;
changing a computation option requires a new fit or CV evaluation.

\subsection{Comparing and choosing covariate specifications}\label{sec:covariate-comparison}

\code{compare\_imr()} provides structural summaries; prediction-based comparison
uses \code{cv\_imr()} on the same cohort and identical fold assignments.
The companion script demonstrates two prespecified formulas,
\code{y \textasciitilde{} age} and
\code{y \textasciitilde{} age + sex + stage}, on the 300 synthetic subjects.
The covariate names are illustrative; this is not a clinical selection result.
Both candidates use the same assays, priors and 2000 retained plus 500 burn-in
draws. Two rounds of three-fold CV yield mean MSEs of 0.349793 and 0.376508,
respectively. The expanded-minus-age-only differences are 0.020337 and 0.033092.
These paired round summaries are descriptive; overlapping training samples do
not provide independent replicates for a significance test.

For data-driven formula choice, the script additionally performs three-fold
outer validation. Within each outer training set, three-fold inner CV chooses
between the two formulas using MSE, with matched inner partitions. The selected
training fit predicts only its outer test subjects. The script checks the actual
ID/fold assignments and verifies that inner records contain no outer test IDs.
The age-only candidate was selected in all three outer folds; the pooled outer
MSE of this selection procedure was 0.340617. This is a separate evaluation of
the selection procedure, not the smaller of the two preceding CV estimates.

From the extracted replication directory, the following commands print the
comparison and nested-selection results. The main replication script invokes
the same function and retains its tables and subject-level records.
\Needspace{48\baselineskip}
\begin{CodeInput}
R> source("IntegMultiReg-covariate-comparison.R")
R> covariate_comparison <- run_covariate_comparison()
R> print(covariate_comparison$paired_summary, digits = 6)
R> print(covariate_comparison$nested_summary, digits = 6)
\end{CodeInput}
\begin{CodeOutput}

Prespecified candidates: paired CV (MSE; lower is better)
           fit    outcome method platforms subgroups retained_draws
      age_only continuous    imr         3         4           2000
 age_sex_stage continuous    imr         3         4           2000
 selected_features
                 6
                 6
     candidate               formula  mean_mse sd_across_rounds
      age_only               y ~ age 0.3497932      0.012413876
 age_sex_stage y ~ age + sex + stage 0.3765077      0.003394637

Paired round differences (not independent inferential replicates)
 round  age_only age_sex_stage expanded_minus_age_only
     1 0.3585712     0.3789081              0.02033692
     2 0.3410153     0.3741073              0.03309205

Nested selection: the outer responses are used only for scoring
 outer_fold selected n_train n_test inner_mse_age_only
          1 age_only     200    100          0.3648101
          2 age_only     200    100          0.4379168
          3 age_only     200    100          0.4123364
 inner_mse_age_sex_stage outer_mse
               0.3971183 0.4648348
               0.4925576 0.2626127
               0.4597745 0.2944039
                  procedure n_subjects outer_folds inner_folds
 inner-CV formula selection        300           3           3
 pooled_outer_mse
        0.3406171
      candidate               formula mean_mse sd_across_rounds
1      age_only               y ~ age 0.349793       0.01241388
2 age_sex_stage y ~ age + sex + stage 0.376508       0.00339464
                   procedure n_subjects outer_folds inner_folds
1 inner-CV formula selection        300           3           3
  pooled_outer_mse
1         0.340617
\end{CodeOutput}

This workflow selects a predictive specification from a fixed candidate list;
it does not introduce a new clinical-coefficient selection prior. Variables
required for adjustment must remain in every candidate formula. Candidate lists,
feature screening and tuning must be confined to training data if selected
adaptively. The example's narrow synthetic setting does not establish clinical
utility, high-dimensional calibration or MCMC convergence.

\Needspace{5\baselineskip}
\subsection{Computational scale}\label{sec:computational-scale}

The computational cost is driven primarily by the number of retained
availability subgroups, the number of candidate molecular features in the
platforms appearing in those subgroups, and the total number of MCMC iterations.
Refit cross-validation repeats preprocessing and the full MCMC fit for each fold
and round, multiplying the main fitting cost by the number of folds and rounds.
The argument \code{max\_models} controls a second
cost in prediction: the maximum number of distinct sampled selection
models used for Bayesian model averaging.

Every retained selection state remains available in its original order.
Identical platform-specific selection matrices share storage in a newly fitted
object, while ordinary R edits retain copy-on-modify behaviour. The research
runner uses lossless compact checkpoints to preserve this sharing on recovery;
ordinary RDS files preserve the values but may use more memory after reload.
Memory requirements still depend on the number of distinct states and native
workspaces, so full-scale runs require measured resource checks.

Table~\ref{tab:scale} summarises the scale of the workflows used in this
paper.  The simulated-data workflow is designed to be a routine replication
check.  The reduced KIRC workflow exercises the real multi-platform data
structure while remaining suitable for examples. The covariate-comparison
workflow adds paired cross-validation and nested formula selection; its repeated
training fits make it more costly than a single illustrative fit.

\begin{table}[!htbp]
\centering
\small
\begin{tabular}{@{}>{\raggedright\arraybackslash}p{3.2cm}
                r r r
                >{\raggedright\arraybackslash}p{3.9cm}@{}}
\toprule
Workflow & Subjects & Features & Subgroups & Manuscript configuration \\
\midrule
\code{simIMR} figures and Table~\ref{tab:compare} & 300 & 38 & 4 &
Figures use 8000 retained draws and 2000 burn-in; prediction comparison uses
4000 retained draws, 1000 burn-in and 10 rounds of 5-fold cross-validation. \\
\code{kircIMR} reduced example & 521 & 130 & 4 &
Runnable real-data example with 4000 retained draws and 1000 burn-in in the
manuscript code. \\
Covariate comparison and nested selection & 300 & 38 & 4 &
2000 retained draws and 500 burn-in; paired 2 rounds of 3-fold CV, followed by
3 outer folds with 3-fold inner selection. \\
\bottomrule
\end{tabular}
\caption{\label{tab:scale}Computational scale of the main workflows.  Candidate
features exclude clinical covariates.}
\end{table}

The current replication materials generate the simulation results, the
covariate-comparison workflow and the reduced KIRC interface example. Legacy
raw-time KIRC results and their provenance are retained in a separate historical
audit directory; they are not performance evidence for the current package.


\FloatBarrier

\Needspace{5\baselineskip}
\subsection{Software validation}\label{sec:software-validation}

The package includes automated checks at both the \proglang{R} interface and
compiled-code boundary.  The \proglang{R} layer validates identifiers, numeric
feature columns, binary/survival outcome coding, subgroup-size choices,
hyperparameter lengths and prediction inputs before calling the sampler.  This
is especially important for multi-platform data, where a silent row-order
mismatch between outcome, covariate and platform data would invalidate the
analysis.  The compiled routines follow standard \proglang{R} registration and
error-handling practices so that numerical failures are reported through the
\proglang{R} interface.

The test suite is organised around the public workflow, covering input
validation, subgroup construction, helper functions, model fitting,
reproducibility under fixed seeds, continuous, binary and survival outcome handling,
S3 methods, plots, prediction, cross-validation and
the KIRC data object.
Local package checks verify installation, examples, tests, vignette rebuilding,
compiled-code checks and manual building.
The revision adds independent checks for computational conventions. Frozen
package results cover six outcome/model configurations and eighteen CV
combinations. An independently compiled function from the 2017 archive checks
conditional CV predictions at fixed states, latent responses and ordered
partitions. Finite-state transition checks establish detailed balance for the
paper selection update, and native-chain frequencies are compared with
exhaustively enumerated conditional targets. Gamma log-density increments are
checked against \code{dgamma()}, including interactions below 0.001.

Historical regression tests remain in place for response scaling, formula
transformations, binary probability averaging, native garbage-collection
protection, subgroup-pair labels, worker consistency and numerical-failure
cleanup. Source-package checks and instrumented CI have separate roles;
previous remote sanitizer or Valgrind results are not evidence that this
unpublished candidate has passed those jobs. The accompanying validation
report identifies the source, environment, completed checks and outstanding
experiments. Conditional coefficient intervals retain the selection sampler's
Laplace approximation and finite-chain uncertainty.
The strongest external check is an original-scale reproduction of the published
application. Running the released 778/91/729-feature matrices under
\code{reference = "code2017"}, with the article's 350000 retained draws after
50000 burn-in and its ten rounds of ten-fold cross-validation, returns the
C-indices in Table~\ref{tab:kirc-reproduction}. The largest absolute difference is
0.044 and the median is 0.011. Seventeen of the twenty-five cells fall inside
the standard error \citet{chekouo2017} report for that cell and twenty-four
fall inside twice it; the exception is the pooled column of IMR $+$ C $+$ M,
where the difference is 0.020 against a reported standard error of 0.01. The archive does not distribute the replicate seeds, so the
agreement is within sampling variation rather than at the level of digits.

\begin{table}[!ht]
\centering
\begin{tabular}{l c c c c c}
\toprule
 & \multicolumn{5}{c}{Availability subgroup} \\
\cmidrule(lr){2-6}
Model & $E_1$ (\code{111}) & $E_2$ (\code{011}) & $E_3$ (\code{101}) &
$E_5$ (\code{001}) & Full \\
 & $n=128$ & $n=70$ & $n=120$ & $n=130$ & $n=448$ \\
\midrule
CPH $+$ C & 0.72 (0.01) & 0.76 (0.01) & 0.79 (0.01) & 0.71 (0.01) & 0.71 (0.01) \\
 & \textit{0.733} (0.014) & \textit{0.757} (0.018) & \textit{0.784} (0.011) & \textit{0.702} (0.016) & \textit{0.721} (0.007) \\[2pt]
IMR $+$ C $+$ M & 0.89 (0.02) & 0.82 (0.04) & 0.89 (0.02) & 0.80 (0.04) & 0.86 (0.01) \\
 & \textit{0.859} (0.009) & \textit{0.829} (0.016) & \textit{0.879} (0.009) & \textit{0.775} (0.013) & \textit{0.840} (0.004) \\[2pt]
BMS $+$ C $+$ M & 0.87 (0.03) & 0.82 (0.04) & 0.89 (0.02) & 0.79 (0.04) & 0.85 (0.01) \\
 & \textit{0.885} (0.009) & \textit{0.864} (0.014) & \textit{0.887} (0.006) & \textit{0.798} (0.014) & \textit{0.859} (0.004) \\[2pt]
IMR $+$ M & 0.88 (0.02) & 0.71 (0.06) & 0.84 (0.02) & 0.72 (0.05) & 0.82 (0.01) \\
 & \textit{0.910} (0.007) & \textit{0.674} (0.023) & \textit{0.865} (0.008) & \textit{0.688} (0.016) & \textit{0.823} (0.004) \\[2pt]
BMS $+$ M & 0.88 (0.02) & 0.67 (0.06) & 0.84 (0.02) & 0.72 (0.05) & 0.81 (0.01) \\
 & \textit{0.849} (0.010) & \textit{0.693} (0.024) & \textit{0.851} (0.009) & \textit{0.719} (0.017) & \textit{0.808} (0.005) \\
\bottomrule
\end{tabular}
\caption{\label{tab:kirc-reproduction}Original-scale KIRC predictive
performance. Upright entries are the published C-indices of
\citet[Table~1]{chekouo2017}; italic entries are the rerun produced by
\pkg{IntegMultiReg} under \code{reference = "code2017"} with 350000 retained
draws after 50000 burn-in, ten rounds of ten-fold cross-validation and the mean
taken over the 100 validation sets. Parenthetic values are standard errors;
the published standard errors are reported to two decimals. Subgroup labels
follow the bitstring convention of Table~\ref{tab:bitstring-key}, where the
right-most digit denotes mRNA. The largest absolute difference is 0.044 and the median
is 0.011; seventeen cells fall inside the published standard error for that
cell and twenty-four inside twice it.
Historical replicate seeds are unavailable, so this is agreement within
sampling variation rather than digit-level reproduction.}
\end{table}


The simulation study reproduces as well, and at the level of its mechanism
rather than only its numbers. Table~\ref{tab:simulation-reproduction} runs the
archived data generator for the three overlap scenarios at 50 replicates each
and the six correlation settings at 30 replicates each, the counts the article
specifies. IMR exceeds BMS on selection AUC in every platform-by-subgroup cell
when the relevant biomarker sets coincide, in every cell at fifty percent
overlap by a smaller margin, and in two of eight cells by at most 0.005 when
the sets do not overlap. That ordering is what the MRF prior predicts and what
\citet{chekouo2017} report; pooled predictive C-indices differ by at most 0.004
between the two methods in all three scenarios, matching the article's
description of IMR as competitive on prediction while better on selection.
Selection AUC for IMR stays near 0.90 as the marker correlation rises to 0.8.

\begin{table}[!ht]
\centering
\begin{tabular}{l l c c c c}
\toprule
\multicolumn{6}{l}{\textit{(a) Main simulation, 50 replicates per scenario}} \\
\midrule
Scenario & Biomarker overlap & Cells favouring & IMR $-$ BMS &
\multicolumn{2}{c}{Mean selection AUC} \\
\cmidrule(lr){5-6}
 & & IMR & range & IMR & BMS \\
\midrule
1 & identical sets   & 8/8 & $+0.022$ to $+0.128$ & 0.831 & 0.757 \\
2 & 50 percent       & 8/8 & $+0.007$ to $+0.062$ & 0.798 & 0.767 \\
3 & none             & 2/8 & $-0.027$ to $+0.005$ & 0.762 & 0.771 \\
\midrule
\multicolumn{6}{l}{\textit{(b) Correlated-marker simulation, 30 replicates per setting}} \\
\midrule
Correlated & Correlation & \multicolumn{2}{c}{} &
\multicolumn{2}{c}{Mean selection AUC (SD)} \\
\cmidrule(lr){5-6}
markers & $\rho$ & & & IMR & BMS \\
\midrule
all  & 0.2 & & & 0.898 (0.055) & 0.745 (0.095) \\
all  & 0.5 & & & 0.891 (0.057) & 0.738 (0.098) \\
all  & 0.8 & & & 0.840 (0.074) & 0.751 (0.093) \\
half & 0.2 & & & 0.916 (0.050) & 0.776 (0.093) \\
half & 0.5 & & & 0.943 (0.038) & 0.840 (0.077) \\
half & 0.8 & & & 0.956 (0.036) & 0.874 (0.067) \\
\bottomrule
\end{tabular}
\caption{\label{tab:simulation-reproduction}Original-scale reproduction of the
simulation study of \citet[Section~5.2 and Web Appendix~E]{chekouo2017}, run
under \code{reference = "code2017"}. Cells are the eight platform-by-subgroup
combinations in which a selection AUC is defined; a panel entry averages over
those cells. Panel (a) shows that the IMR advantage over BMS is large when the
relevant biomarker sets coincide, smaller at fifty percent overlap, and absent
when they do not overlap, which is the ordering the article predicts from the
MRF prior. Pooled predictive C-indices are close throughout (IMR 0.819, 0.822,
0.836 against BMS 0.815, 0.819, 0.835), matching the article's description of
IMR as competitive on prediction while better on selection. Panel (b) keeps the
six settings of Web Appendix~E separate rather than pooling them: the first
three correlate every signal marker with the unobserved gene and the last three
correlate half of them. Parenthetic values are the between-replicate standard
deviation, averaged over cells. IMR stays above BMS in every setting, but the
direction in $\rho$ differs between the two halves of the design, so the two
should not be averaged together.}
\end{table}


These checks are not a substitute for substantive model diagnostics, but they
reduce the chance that a package user reaches the MCMC stage with misaligned
data or an unsupported prediction request.

%% ===========================================================================
\section{Illustration}\label{sec:illustration}
%% ===========================================================================

This section evaluates selection and prediction on simulated data with known
truth. Appendix~\ref{sec:appendix-kirc} reports sample retention and provides
a runnable illustration using the reduced KIRC data.

The package ships with two example data sets.  The first, \code{kircIMR}, is a
reduced public UCSC Xena TCGA-KIRC data set that follows the kidney cancer
case-study structure of \citet{chekouo2017}.  It is built from public sampleMap
files, not controlled-access TCGA/GDC files, and stores a reduced screened
feature panel with package-internal patient identifiers.  The second,
\code{simIMR}, is a simulated data set with known truth; we use it for the
quantitative recovery and prediction summaries because the true biomarkers are
known and the full analysis is fast to reproduce.

\subsection{Simulated example with known truth}\label{sec:sim-illustration}

We use the simulated \code{simIMR} data, in which a handful of features per
platform (recorded in \code{simIMR\$truth}: \code{G01}--\code{G03},
\code{P01}--\code{P02}, \code{M01}) drive the outcomes through a common latent
signal. The four non-empty subgroups are \code{011} ($n=120$, genomic +
proteomic), \code{111} ($n=60$, all three), \code{101} ($n=60$, genomic +
metabolomic) and \code{100} ($n=60$, metabolomic only). All numbers and figures
below are reproduced by the accompanying replication scripts.

\subsection{Biomarker recovery}\label{sec:recovery}

Table~\ref{tab:recovery} contrasts the largest inclusion probability among
associated and null features on each platform. Each platform contains at least one true feature whose largest subgroup mPIP
rounds to 1.000, while the corresponding null-feature maxima range from 0.178
to 0.712. These maxima describe the strongest selection signals; they do not
measure recovery of every true feature or the false-positive rate. The high
null-feature probabilities also motivate checking sensitivity to sparsity and
MRF borrowing before interpreting the selected set.

\begin{table}[!htbp]
\centering
\begin{tabular}{l r r r r}
\toprule
Platform & \#features & \#true & max mPIP (true) & max mPIP (null) \\
\midrule
genomic     & 20 & 3 & 1.000 & 0.712 \\
proteomic   & 10 & 2 & 1.000 & 0.178 \\
metabolomic &  8 & 1 & 1.000 & 0.505 \\
\bottomrule
\end{tabular}
\caption{\label{tab:recovery}Largest marginal inclusion probability among the
truly associated versus the null features of each platform (right-censored
outcome, 8000 retained MCMC draws).}
\end{table}

The table reports
the maximum mPIP over availability subgroups, matching the ranking used by
\code{summary()} and \code{plot\_top\_features()}.

\subsection{Comparing integration with independent subgroup models}\label{sec:gain}

Borrowing strength can change both selection and prediction. Its effect is
evaluated against the no-borrowing BMS comparator. Using the training-fold refitting cross-validation
assessment described in Section~\ref{sec:predict},
Table~\ref{tab:compare} compares the integrative model (\code{method = "imr"})
with the non-integrative variant (\code{method = "bms"}) that fits every
subgroup independently. For each of the three outcome types we report the
out-of-fold predictive accuracy
($5$ folds, averaged over $10$ rounds) per subgroup and overall.

\begin{table}[t!]
\centering
\begin{tabular}{l l l r r r r r}
\toprule
Outcome & Measure & Method & \code{011} & \code{100} & \code{101} & \code{111} & all \\
\midrule
binary & AUC ($\uparrow$) & IMR & \textbf{0.660} & \textbf{0.640} & 0.605 & \textbf{0.605} & \textbf{0.628} \\
 &  & BMS & 0.609 & 0.617 & \textbf{0.674} & 0.504 & 0.607 \\
continuous & MSE ($\downarrow$) & IMR & 0.298 & 0.271 & \textbf{0.375} & \textbf{0.479} & \textbf{0.344} \\
 &  & BMS & \textbf{0.294} & \textbf{0.270} & 0.492 & 0.658 & 0.402 \\
right-censored & C-index ($\uparrow$) & IMR & 0.854 & 0.633 & 0.762 & \textbf{0.811} & \textbf{0.801} \\
 &  & BMS & \textbf{0.861} & \textbf{0.650} & \textbf{0.764} & 0.761 & 0.794 \\
\bottomrule
\end{tabular}
\caption{\label{tab:compare}Cross-validated predictive accuracy of the
integrative (IMR) and non-integrative (BMS) models by subgroup and overall,
for each outcome type ($5$-fold cross-validation averaged over $10$ rounds).
The assessment uses independent training-fold refits described in
Section~\ref{sec:predict}. Arrows indicate the favourable direction; the better
of the two point estimates is in \textbf{bold}, without implying statistical
significance. The overall column pools subject predictions within each round
before averaging the scores across rounds.}
\end{table}

In this example, overall IMR AUC is 0.628 versus 0.607 for BMS, MSE is
0.344 versus 0.402, and concordance is 0.801 versus 0.794. The subgroup
comparisons are mixed: IMR has lower continuous MSE in \code{101} and
\code{111}, but higher MSE in \code{011} and \code{100}. Survival
concordance is higher under BMS in three of the four subgroups. These are
results for the explicit paper sampler and fixed priors on this example;
they do not establish a general predictive advantage or assess Monte Carlo
uncertainty across independent data sets.


\subsection{Diagnostics and sensitivity checks}\label{sec:diagnostics}

The examples above use fixed seeds so that the printed tables and figures are
exactly reproducible.  In a final applied analysis, we recommend treating the
reported fit as one element of a small diagnostic workflow.  First, the
log-posterior trace should be inspected for gross non-stationarity or numerical
instability; \code{plot(fit, type = "trace")} provides this check directly from
the fitted object.  Second, the leading selected features should be compared
across several seeds or longer chains.  In high-dimensional biomarker
selection, exact mPIPs for borderline features may change across runs, but the
top-ranked features and the separation between strong signals and null features
should be reasonably stable.

The package makes that second check a routine operation. The archived 2017
program fixes its random seed in the source and exposes no seed argument, so the
eight-chain multi-start assessment reported by \citet{chekouo2017} cannot be
obtained from it without editing and recompiling the \proglang{C} sources. Here
\code{initial} takes explicit selection and interaction starting values and
\code{seed} is an ordinary argument, so a multi-start run is scripted rather
than patched. Table~\ref{tab:chain-diagnostics} reports eight chains from eight
recorded starts, each retaining $10^{6}$ draws after 50000 burn-in, roughly
three times the per-chain budget of the published analysis. The mRNA and miRNA
interaction parameters reach rank-normalized split-$\widehat{R}$ below 1.01
with bulk effective sample sizes in the thousands; the methylation interaction
does not, at $\widehat{R}=1.045$ and bulk ESS 208. The smallest cross-chain
correlation between marginal posterior inclusion probabilities is 0.803 for
mRNA and 0.813 for methylation, below the 0.85 bound the article reports, while
miRNA reaches 0.958.

\begin{table}[!ht]
\centering
\begin{tabular}{l c c c c c c}
\toprule
Platform & Pairs & Rank $\widehat{R}$ & Bulk ESS & Tail ESS & Classical PSRF &
Min.\ $\rho$ \\
\midrule
mRNA        & 6 & 1.001--1.001 & 7,832--15,997 & 26,775--49,970 & 1.000--1.001 & 0.803 \\
miRNA       & 1 & 1.007        & 1,426         & 1,490          & 1.005        & 0.958 \\
Methylation & 1 & 1.045        & 208           & 498            & 1.016        & 0.813 \\
\bottomrule
\end{tabular}
\caption{\label{tab:chain-diagnostics}Eight-chain diagnostics for the
subgroup-interaction parameters $\theta$, from eight recorded starting points
with $10^6$ retained draws per chain after 50000 burn-in under
\code{sampler\_method = "paper"}. ``Pairs'' counts the subgroup pairs a
platform contributes; ranges run over those pairs. Rank $\widehat{R}$ is the
rank-normalized split-folded statistic and PSRF the classical potential scale
reduction factor. Min.\ $\rho$ is the smallest Pearson correlation between
marginal posterior inclusion probabilities over all 28 chain pairs and
availability subgroups of that platform.}
\end{table}


The two runs target different stationary distributions: the released selection
update omits a Hastings ratio at boundary flips, and the corrected
\code{"paper"} update used here does not (Section~\ref{sec:strategies}). The
archive records neither the per-chain budget behind the published bound nor the
form of its correlation summary. For an applied analysis the consequence is
that methylation rankings from a single chain of this length are not stable.

The comparison between IMR and BMS is also a sensitivity analysis for the MRF
borrowing assumption.  Agreement between the fits indicates limited sensitivity to the borrowing
choice in that comparison; it does not identify where the evidence originates.
Differences in mPIPs or prediction show that the choice matters, but their
direction need not favour IMR. Shared feature effects across availability
patterns provide a scientific rationale for borrowing, which should be assessed
against subject-matter knowledge as well as the no-borrowing fit.  Similarly, changing \code{nu} probes sensitivity to
the prior sparsity level, and changing the reporting threshold in
\code{summary(fit, threshold = ...)} distinguishes robust high-probability
markers from features near the conventional 0.5 cutoff.

For prediction, subgroup-specific performance is as important as the pooled
column.  A model can have good overall AUC or C-index while performing poorly in
a small availability subgroup.  The column-labelled output of \code{cv\_imr()}
is therefore designed to make this failure mode visible.  The simulated comparison contains both improvements and declines by subgroup,
illustrating why the pooled score alone is insufficient.

%% ===========================================================================
\section{Reproducibility and replication materials}\label{sec:reproducibility}
%% ===========================================================================

The current manuscript's numerical results are generated by
\code{IntegMultiReg-replication.R} and its covariate-comparison companion.
Figures and biomarker recovery use 8000 retained and 2000 burn-in draws;
prediction comparisons use 4000 retained and 1000 burn-in draws with ten rounds
of five-fold training-refit validation. The covariate example uses its separately
reported paired and nested validation settings. Fixed seeds, subject-level fold
assignments, numerical outputs and session information are retained. Quick mode
is a smoke test, not a source of the manuscript numbers.

The synthetic \code{simIMR} data provide known feature truth and the predictive
examples. The reduced public UCSC Xena \code{kircIMR} data exercise the real
missing-platform structure and the survival interface. Their globally screened
feature panel makes them suitable for an interface demonstration, not an unbiased
estimate of generalization performance. Full-scale KIRC performance is reported
separately, from the original-scale runner rather than from these illustration
data, in Table~\ref{tab:kirc-reproduction}.

The original-scale runner is separate from the interface illustration. It
uses the released 778/91/729-feature matrices for Table 1 and compiles the
archived data generator for the three overlap scenarios and the six
correlation settings of Web Appendix E. A full configuration retains 350000
draws after 50000 burn-in iterations; Table 1 requests ten rounds of ten-fold
CV and reports the mean over 100 validation sets, with SD and standard error
stored separately. The simulation design distinguishes released fixed marker indices from
random marker identities described in the article. A common seeded permutation
per platform maps the archived signal pattern to random identities while
preserving signal counts and overlap; output columns and truth labels are
restored to the original feature order. Actual permutations and inner L1-CPH
fold assignments are saved. Unknown historical replicate seeds and the article's
776-versus-archive's-778 gene count remain documented.
These differences prevent an automatic claim of historical-digit reproduction.
The \code{initial} argument supports explicitly different selection and
interaction starts for multi-chain diagnostics. Historical starting values are
unavailable, so the eight recorded starts of Table~\ref{tab:chain-diagnostics}
are ours rather than the article's, and they run the corrected \code{"paper"}
update rather than the released one. Those chains are full-scale, at $10^{6}$
retained draws each, but they characterize the stated model and do not
reproduce the published Appendix F diagnostics.
The archive's global screening is retained as an experimental condition;
training-fold screening is a separate predictive evaluation choice.

Historical scripts and previous numerical references are retained in
\code{historical-audit}. Short runs verify orchestration and are kept separate
from full-scale scientific results. The reproduction status, elapsed times,
source/data hashes and outstanding analyses are recorded in the accompanying
report. No proprietary IPA network analysis is regenerated.

%% ===========================================================================
\section{Summary}\label{sec:summary}
%% ===========================================================================

\pkg{IntegMultiReg} makes the availability-subgroup model of
\citet{chekouo2017} available through a common \proglang{R} workflow for
survival, continuous and binary outcomes. The data and model interfaces allow
partially observed subjects to contribute through their measured platforms,
subject to the outcome, covariate and subgroup-size requirements. Standard
methods then connect fitting to selection summaries, uncertainty assessment,
prediction and cross-validation.

The simulated illustration shows why both feature-level and subgroup-level
results matter. High inclusion probabilities occur among true and null
features, and the predictive comparisons favour different methods in different
subgroups. The package therefore supports examining the consequences of MRF
borrowing rather than assuming that integration improves every analysis.
Multi-chain diagnostics, prior sensitivity and validation appropriate to the
prediction target remain necessary for substantive applications.

The current model uses platform-specific sparsity parameters $\nu_k$.
Extensions could incorporate known biological interactions among features or
adapt the feature-ordering and population-structure construction of
\citet{dai2023} to this availability-subgroup setting.

%% ---------------------------------------------------------------------------
\section*{Computational details}
The revised simulation and reduced KIRC results were obtained using
\proglang{R}~4.4.2 on macOS~26.3.1 (ARM64), with \pkg{IntegMultiReg}~0.2.0
and GSL~2.8. The supplementary materials include the source package and
workflows reproducing the current figures and simulated-data tables.
Historical audit materials are supplied separately.

\bibliography{IntegMultiReg}

%% ===========================================================================
\appendix
\renewcommand{\thetable}{\thesection.\arabic{table}}
\renewcommand{\thefigure}{\thesection.\arabic{figure}}
\renewcommand{\theequation}{\thesection.\arabic{equation}}
\renewcommand{\theHtable}{appendix.\thesection.\arabic{table}}
\renewcommand{\theHfigure}{appendix.\thesection.\arabic{figure}}
\setcounter{table}{0}
\setcounter{figure}{0}
\setcounter{equation}{0}
\clearpage
\section{Supplementary empirical analyses and outcome extensions}
\label{sec:appendix-kirc}
%% ===========================================================================

\subsection{Complete-case sample retention}\label{sec:complete-case}

The practical motivation for IMR is illustrated by comparing the number of
subjects retained under a complete-case analysis with the number retained by
the subgroup-preserving IMR workflow.  Table~\ref{tab:complete-case} reports
this comparison for the two example data objects.  In both cases the
complete-case analysis keeps only the \code{111} subgroup, i.e., subjects with
all three platforms observed.  IMR instead keeps every availability subgroup
whose size exceeds the modelling threshold.

\begin{table}[!htbp]
\centering
\small
\begin{tabular}{@{}>{\raggedright\arraybackslash}p{4.3cm} r r r r@{}}
\toprule
Data set & Any platform & Complete cases & IMR modelled &
Complete-case \% \\
\midrule
\code{simIMR} & 300 & 60 & 300 & 20.0\% \\
\code{kircIMR} reduced example & 522 & 172 & 521 & 33.0\% \\
\bottomrule
\end{tabular}
\caption{\label{tab:complete-case}Sample retained by complete-case analysis
compared with the subgroup-preserving IMR workflow.  The reduced KIRC example
contains one methylation-only subject below the default \code{min\_subgroup\_size = 30}
threshold, so 521 of 522 subjects are modelled.}
\end{table}

The table also clarifies the role of BMS in the package.  BMS is not a
complete-case method: it retains the same modelled availability subgroups as
IMR, but fits them independently.  Thus the comparison between IMR and BMS
isolates the effect of MRF borrowing, whereas the comparison with complete-case
analysis isolates the effect of retaining partially observed subjects.

\subsection{Reduced package example}\label{sec:kirc}

The original Biometrics analysis integrated mRNA expression, miRNA expression
and DNA methylation with clinical covariates for right-censored survival.  After
pre-processing, the published article reports 776 mRNA transcripts, 91 miRNA
markers and 729 DNA methylation probes \citep{chekouo2017}.  The supplementary
\proglang{C} data files distributed with that article, which are used for the
separate historical audit, contain 778 named
mRNA columns, 91 miRNA columns and 729 methylation columns.  We therefore use
the published screening counts when describing the scientific study and the
supplement matrix dimensions when documenting the source files.  For the
package example, \code{kircIMR} keeps the same scientific structure but stores
a smaller Cox-screened panel (50, 30 and 50 features, respectively), so that the
data object is small enough to ship with the package while still exercising the
multi-platform missingness pattern.

The reduced example is intended to demonstrate the software workflow rather
than to serve as a literal repetition of the original KIRC marker-selection
analysis.  It uses the same three platform types, the same survival-outcome
structure and the same adjustment covariates, but it replaces the long
full-data MCMC analysis by a compact data object suitable for examples, tests
and vignettes.

These distinctions are important because the three KIRC-related objects serve
different replication targets rather than repeated runs of a single identical
analysis: (i) the distributed \code{kircIMR} object is a reduced public example
for users and tests; (ii) a public UCSC Xena reprocessing with approximately the
published feature counts is useful for auditing the data-preparation pipeline;
and (iii) the Biometrics supplementary-data workflow is preserved only in the
separate historical audit.  Differences in
feature counts, patient identifiers, subgroup sizes and selected marker lists
should therefore be interpreted through these data provenance and distribution
boundaries rather than treated as numerical disagreement from a single
repeated analysis.

The shipped object was generated from public UCSC Xena TCGA-KIRC sampleMap
files \citep{tcgaKirc2013,gdcPortal,goldman2020}.  It does not distribute
controlled-access TCGA/GDC data or the full Biometrics/Wiley supplementary data
files.  To reduce unnecessary linkage risk in a package example, the TCGA
barcodes are replaced by package-internal labels such as \code{KIRC001}; no
barcode mapping is distributed, and users should not attempt participant
re-identification or linkage to external resources.

\begin{CodeChunk}
\begin{CodeInput}
R> data("kircIMR")
R> dims <- sapply(kircIMR$platforms, dim)
R> rownames(dims) <- c("subjects", "variables")
R> dims
\end{CodeInput}
\begin{CodeOutput}
          mrna mirna methylation
subjects   521   235         312
variables   51    31          51
\end{CodeOutput}
\begin{CodeInput}
R> kircIMR$model_subgroup_sizes
\end{CodeInput}
\begin{CodeOutput}
bitstrings
011 101 001 111 
 63 139 147 172 
\end{CodeOutput}
\end{CodeChunk}

The subgroup bitstrings use the same convention as the fitted object: with
platform order \code{mrna}, \code{mirna}, \code{methylation}, the right-most bit
corresponds to mRNA and the left-most bit to methylation.  Thus
\code{001} denotes mRNA only, \code{011} mRNA plus miRNA, \code{101} mRNA plus
methylation and \code{111} all three platforms.  These are the subgroups
modelled with the default illustration threshold \code{min\_subgroup\_size = 30}; a single
methylation-only subject is present in the raw availability table but is too
small to model.

A survival analysis on this data set uses the same hyperparameters as the
reference paper for the non-local prior scale and sparsity pattern:

\Needspace{39\baselineskip}
\begin{CodeChunk}
\begin{CodeInput}
R> fit_kirc <- imr(
+    x = kircIMR$platforms,
+    outcome            = kircIMR$outcome.survival,
+    covariates         = kircIMR$covariates,
+    outcome_type       = "right.censored",
+    nu                 = c(-4, -3, -4),
+    draws = 4000, burnin = 1000,
+    sampler_method = "paper",
+    min_subgroup_size              = 30,
+    seed               = 1)
R> summary(fit_kirc, threshold = 0.5)
\end{CodeInput}
\begin{CodeOutput}
Integrative Bayesian Multi-Platform Regression (IMR) -- summary
--------------------------------------------------------------
Outcome type : right.censored   Method: IMR
Selection threshold (mPIP) : 0.50

Platform 'mrna': 3 selected feature(s)
 feature max_mpip subgroup
    PRH2    1.000      111
   NOP56    0.731      001
  COPS7B    0.525      011

Platform 'mirna': 8 selected feature(s)
      feature max_mpip subgroup
 MIMAT0000450    1.000      111
 MIMAT0017994    1.000      111
 MIMAT0004671    1.000      111
 MIMAT0004588    0.986      111
 MIMAT0004556    0.953      111
 MIMAT0003880    0.934      111
 MIMAT0000099    0.560      011
 MIMAT0000076    0.516      011

Platform 'methylation': 0 selected feature(s)

\end{CodeOutput}
\end{CodeChunk}

This compact fit is a runnable real-data example aligned with the case study.

\subsection{Continuous-outcome extension}\label{sec:appendix-continuous}

For a retained subgroup $s$, write the linear predictor as
%
\begin{equation}\label{eq:eta-continuous}
  \bm{\eta}^{(s)}
  = \beta_0^{(s)}\bm{1}_{n_s}
  + \bm{C}^{(s)}\bm{\beta}_{C}^{(s)}
  + \bm{X}_{\gamma}^{(s)}\bm{\beta}_{\gamma}^{(s)},
\end{equation}
%
where $\bm{X}_{\gamma}^{(s)}$ contains only the molecular features selected by
the current selection model $\bm{\gamma}$, and
$\bm{\beta}_{\gamma}^{(s)}$ denotes the corresponding nonzero regression
coefficients.  For a continuous response no augmentation is required and
%
\begin{equation}\label{eq:continuous-model}
  \bm{y}^{(s)}=\bm{y}^{\star(s)}, \qquad
  \bm{y}^{(s)}\mid\bm{\eta}^{(s)},\sigma_s^2
  \sim \mathcal{N}\!\left(\bm{\eta}^{(s)},\sigma_s^2\bm{I}_{n_s}\right).
\end{equation}
%
The corresponding likelihood is therefore
%
\begin{equation}\label{eq:continuous-likelihood}
  L_s(\bm{\beta}_{\gamma}^{(s)},\sigma_s^2)
  \propto (\sigma_s^2)^{-n_s/2}
  \exp\!\left\{
    -\frac{1}{2\sigma_s^2}
    \left\|\bm{y}^{(s)}-\bm{\eta}^{(s)}\right\|^2
  \right\}.
\end{equation}
%
Combining this Gaussian likelihood with the inverse-gamma prior on
$\sigma_s^2$ and pMOM priors on all active coefficients, including the
intercept and clinical effects, gives the same marginal
likelihood calculation used by the sampler.  Thus the continuous extension is
obtained by taking the observed response itself as the latent Gaussian
response.

\subsection{Binary-outcome extension}\label{sec:appendix-binary}

For a binary outcome, the package uses the standard latent Gaussian probit
data augmentation representation of \citet{albert1993}; the same construction
is used in Bayesian predictive
imaging-genetics models for binary case-control classification
\citep{chekouo2016aoas}.  For subject $i$ in subgroup $s$, let
%
\begin{equation}\label{eq:binary-latent}
  z_i^{(s)} = \eta_i^{(s)} + e_i^{(s)}, \qquad
  e_i^{(s)} \sim \mathcal{N}(0,\sigma_s^2),
\end{equation}
%
where $\eta_i^{(s)}$ is the $i$th element of the linear predictor in
Equation~\ref{eq:eta-continuous}.  The latent variable
$z_i^{(s)}$ is continuous only inside the sampler; the data supplied by the user
remain the observed binary labels.  The observed binary label is generated by
thresholding the latent Gaussian utility,
%
\begin{equation}\label{eq:binary-threshold}
  y_i^{(s)} = I\!\left(z_i^{(s)} > 0\right).
\end{equation}
%
Marginalizing over $z_i^{(s)}$ gives the probit likelihood
%
\begin{equation}\label{eq:binary-probit}
  \Pr\!\left(y_i^{(s)}=1\mid\eta_i^{(s)}\right)
  = \Phi\!\left(\eta_i^{(s)}/\sigma_s\right), \qquad
  \Pr\!\left(y_i^{(s)}=0\mid\eta_i^{(s)}\right)
  = 1-\Phi\!\left(\eta_i^{(s)}/\sigma_s\right),
\end{equation}
%
where $\Phi(\cdot)$ is the standard normal distribution function. The
implementation approximates unit variance using an inverse-gamma prior with
shape and rate $10^5$; it does not impose an exact point mass at one. In the
conditional coefficient sampler, the probability is therefore
$\Phi(\eta_i^{(s)}/\sigma_s)$ for each variance draw. Because the
regression coefficients $\bm{\beta}_{\gamma}^{(s)}$ are integrated out of the
marginal likelihood under the non-local pMOM prior, the latent utilities do not
have a closed-form Gaussian full conditional and are instead updated with a
Metropolis--Hastings step.  Following \citet{chekouo2016aoas}, a candidate
$\tilde{z}_i^{(s)}$ that respects the sign constraint implied by $y_i^{(s)}$ is
drawn from a one-sided exponential proposal with mean $1/|z_i^{(s)}|$,
%
\begin{equation}\label{eq:binary-latent-update}
  \tilde{z}_i^{(s)} \sim \mathrm{Exp}_{\mathrm{mean}}\!\left(1/|z_i^{(s)}|\right)
  \ \text{if } y_i^{(s)}=1,
  \qquad
  \tilde{z}_i^{(s)} \sim -\,\mathrm{Exp}_{\mathrm{mean}}\!\left(1/|z_i^{(s)}|\right)
  \ \text{if } y_i^{(s)}=0,
\end{equation}
%
and accepted with the Metropolis--Hastings probability formed from the non-local
marginal likelihood.  The proposal keeps observed $1$s at positive latent
utilities and observed $0$s at negative ones, so the sign constraint of the
probit augmentation is preserved throughout the chain.
%
Conditional on the augmented vector
$\bm{z}^{(s)}=(z_1^{(s)},\dots,z_{n_s}^{(s)})^\top$, the binary model again has
a Gaussian working likelihood,
%
\begin{equation}\label{eq:binary-working}
  \bm{z}^{(s)}\mid\bm{\eta}^{(s)}
  \sim \mathcal{N}\!\left(\bm{\eta}^{(s)},\sigma_s^2\bm{I}_{n_s}\right).
\end{equation}
%
The variable-selection step can therefore use the same selected design matrix
$\bm{X}_{\gamma}^{(s)}$, the same pMOM prior on the active molecular effects,
and the same MRF prior that links the indicators across availability subgroups.
For point prediction, the probit link is applied within each retained selection
model before Bayesian model averaging. If $w_m$ is the normalized predictive
weight and $\hat{\eta}_m$ the fitted latent linear predictor for model $m$, the
implementation uses the unit-variance approximation
\[
  \widehat{\Pr}(y=1\mid\text{data})
  = \sum_m w_m\,\Phi(\hat{\eta}_m), \qquad \sum_m w_m=1.
\]
Thus, probabilities are averaged on the response scale. Conditional posterior
prediction instead evaluates $\Phi(\eta/\sigma)$ for each sampled coefficient
and variance draw, retaining uncertainty as described in the fitted-object
section.

\end{document}
